%!TEX root = ../main.tex

\addchap{Abstract}

Klassische \ac{ETL}-Prozesse sind regelbasiert und setzen voraus, dass jede
Verarbeitungsvorschrift zuvor vollständig formuliert wird. Die vorliegende
Arbeit untersucht experimentell, inwieweit \acp{LLM} diese Aufgabe übernehmen
können. Auf einem synthetischen Datensatz mit kontrolliert eingebauten
Qualitätsproblemen werden vier Modelle, drei Prompting-Strategien und neun
\ac{ETL}-Aufgaben gegen eine regelbasiert implementierte Vergleichspipeline
gemessen, ergänzt um eine verkettete Ausführung, einen isolierten Durchlauf zur
Trennung von eigener Schrittleistung und geerbtem Fehler, einen Vergleich zweier
Ausführungsformen sowie zwei Ebenen zum Zuschnitt der Aufgabenstellung.
Insgesamt liegen 711 Modellaufrufe zugrunde.

Die Ergebnisqualität hängt weit stärker vom Aufgabentyp ab als vom Modell:
Zwischen dem besten und dem schwächsten kommerziellen Modell liegen rund acht
Prozentpunkte, zwischen der am besten und der am schlechtesten gelösten Aufgabe
nahezu hundert. Als zweite Einflussgröße erweist sich die Ausführungsform. Der
erwartete Vorteil bei semantisch geprägten Aufgaben tritt allein in der direkten
Verarbeitung ein, während die Code-Generierung zuverlässiger rechnet und je
Zeile rund ein Fünfzehntel kostet. Am stärksten wirkt der Zuschnitt der
Aufgabe: Dieselbe Verarbeitung bleibt als eine zusammengesetzte Aufgabe
ungelöst und wird in Teilschritte zerlegt fehlerfrei bewältigt, wobei der
messbare Unterschied zwischen den Prompting-Strategien dabei vollständig
verschwindet. Die Arbeit leitet daraus eine hybride Architektur ab, die
regelbasierte und modellgestützte Schritte je nach Aufgabentyp kombiniert und
nach jedem Schritt inhaltlich prüft.

\vspace{1em}
\noindent
\textit{Classical \ac{ETL} processes are rule-based and require every
transformation rule to be specified in advance. This thesis experimentally
evaluates to what extent \acp{LLM} can take over that task. Using a synthetic
dataset with deliberately injected quality defects, four models, three prompting
strategies and nine \ac{ETL} tasks are measured against a rule-based reference
pipeline, complemented by a chained execution, an isolated run separating a
step's own performance from inherited error, a comparison of two execution
modes, and two levels addressing how the task itself is framed. The evaluation
rests on 711 model calls.}

\vspace{0.6em}
\noindent
\textit{Result quality depends far more on the type of task than on the model:
roughly eight percentage points separate the best from the weakest commercial
model, whereas nearly a hundred separate the best-solved from the worst-solved
task. The execution mode emerges as a second decisive factor. The expected
advantage on semantically demanding tasks materialises only under direct
processing, while code generation computes more reliably at about one fifteenth
of the cost per row. The strongest lever, however, is task framing: the same
transformation remains unsolved when posed as one composite task and is
completed without error once decomposed into individual steps, at which point
the measurable difference between prompting strategies disappears entirely. On
this basis the thesis derives a hybrid architecture that combines rule-based and
model-based steps according to task type and validates the result after every
step.}
